\scanpage{148}
% archive foliation: 147
\begin{center}42\end{center}
let $D'_h=\widetilde\Phi^{-1}(D_h^{*\prime})$. Clearly $D_h^{*\prime}$ meets $D_h^*$ at simple intersection points only and therefore $(D_h^{*2})=(D_h^*D_h^{*\prime})$ is equal to the number of these intersection points, while each intersection point of $D_h^{*\prime}$ and $D_h^*$ corresponds to $m$ intersection points of $D'_h$ and $D_h$. Hence we have
\[
(D_h^2)=(D_hD'_h)\geqq m(D_h^{*2})=2mah^2,
\]
and therefore
\[
\dim|D_h|\geqq mah^2-bh+\chi(\widetilde V)-1-\widetilde p_g,
\]
where $2b=(KD_1)$ and $\widetilde p_g$ is the geometric genus of $\widetilde V$. Combined with the equality
\[
\dim|D_h|=\dim|D_h^*|=ah^2-b^*h+\chi(V*)-1,
\]
this proves that
\[
mah^2-bh\leqq ah^2-b^*h+\text{const.},\quad\text{for all }h\geqq1,
\]
where $a=\frac12(D_1^{*2})$ is positive. Consequently we obtain $m=1$.
Thus the inverse image $\widetilde\Phi^{-1}(\mathfrak p*)$ of the general point $\mathfrak p*$ on $V*$ consists of a single point and therefore $\widetilde\Phi^{-1}$ has no branch curve. Hence $\widetilde\Phi^{-1}$ is regular and single-valued except for a finite number of fundamental points. $\widetilde\Phi$ is therefore bi-regular on $\widetilde V$ except for a finite number of fundamental curves.

Now these fundamental curves can be eliminated by quadratic transformations. Namely, by applying a finite number of

